%----------------------------------------------------------------------------------------
% ONDERLINGE LIGGING. Alle opgaven zijn eigen werk.
%
% Opbouw: Start, leerstof (theorie, werkwijze of uitleg, voorbeeld, korte oefening), en op een
% nieuwe pagina de Opdrachten (door elkaar, stijgende moeilijkheid), Koppel en Test jezelf.
% Labels van dit hoofdstuk beginnen met h3-.
%----------------------------------------------------------------------------------------

\mychapter{3}{Onderlinge ligging}

\lkr{Leerplan B+S'': LPD 5-8 (ruimtemeetkunde), het deel over de onderlinge ligging van rechten en vlakken en het snijden ervan. Kruisende rechten zijn nieuw ten opzichte van het vlak: laat leerlingen ze eerst aanwijzen in het klaslokaal (een rand van het plafond en een rand van de vloer). De stelsels sluiten aan bij Matrices.}

In het vlak snijden twee rechten elkaar, of ze zijn evenwijdig. In de ruimte is er een derde mogelijkheid: kijk naar een rand van het plafond en een rand van de vloer die er dwars op staat. Ze snijden elkaar niet, en toch zijn ze niet evenwijdig. Zulke rechten heten kruisend.

\begin{oefeningen}[Start]
\oef{h3-start-1}

\oef{h3-verken}
\end{oefeningen}

%========================================================================================
\section{Twee rechten}

\naast{%
\begin{theorie}
Twee rechten in de ruimte zijn
\begin{itemize}
\item \textbf{\textcolor{accent}{evenwijdig}} als hun richtingsvectoren collineair zijn (en \textbf{samenvallend} als ze ook een punt gemeen hebben);
\item \textbf{\textcolor{accent}{snijdend}} als ze precies één punt gemeen hebben;
\item \textbf{\textcolor{accent}{kruisend}} als ze niet evenwijdig zijn en geen punt gemeen hebben.
\end{itemize}
\end{theorie}}{%
\begin{tikzpicture}[ruimte,scale=.75,font=\scriptsize]
  \draw[grijs] (2,0,0) -- (2,2,0) -- (0,2,0) (2,0,0) -- (2,0,2) -- (0,0,2) -- (0,2,2) -- (2,2,2) -- (2,0,2) (2,2,0) -- (2,2,2) (0,2,0) -- (0,2,2);
  \draw[grijs,dashed] (0,0,0) -- (2,0,0) (0,0,0) -- (0,2,0) (0,0,0) -- (0,0,2);
  \draw[accent,very thick] (0,0,0) -- (2,0,0);
  \draw[accent,very thick] (0,0,2) -- (0,2,2);
  \node[accent,below left] at (1,0,0) {$l$};
  \node[accent,above] at (0,1,2) {$m$};
\end{tikzpicture}}

\voorbeeld{}{Onderzoek de ligging van $l\colon (1, 0, 2) + r\vcr{1}{1}{0}$ en $m\colon (0, 3, 0) + s\vcr{1}{-1}{1}$.\label{vb:h3-snijdend}}{%
\uitleg{De richtingsvectoren zijn niet collineair. Zoek een gemeenschappelijk punt: stel de coördinaten gelijk.}
\begin{stappenbreed}
1 + r = s,\quad r = 3 - s,\quad 2 = s & Drie vergelijkingen, twee onbekenden. \\
s = 2,\ r = 1,\ \text{controle } 1 + 1 = 2 & De eerste klopt ook. \\
\text{snijpunt } (2, 1, 2) & $r = 1$ invullen in $l$.
\end{stappenbreed}}

\begin{oefeningen}
\oef{h3-snijdend}
\end{oefeningen}

\voorbeeld{}{In de kubus hierboven ligt $l$ op de $x$-as en $m\colon (0, 0, 2) + s\vcr{0}{1}{0}$. Toon aan dat ze kruisend zijn.\label{vb:h3-kruisend}}{%
\uitleg{De richtingsvectoren $\vcr{1}{0}{0}$ en $\vcr{0}{1}{0}$ zijn niet collineair. Zoek een gemeenschappelijk punt.}
\begin{stappenbreed}
r = 0,\quad 0 = s,\quad 0 = 2 & De derde vergelijking kan niet. \\
\text{geen gemeenschappelijk punt} & Niet evenwijdig en niet snijdend: kruisend.
\end{stappenbreed}}

\begin{oefeningen}
\oef{h3-kruisend}
\end{oefeningen}

\begin{werkwijze}
Vergelijk de richtingsvectoren. Collineair: evenwijdig; ligt een punt van de ene op de andere, dan samenvallend. Niet collineair: stel de parametervergelijkingen gelijk (gebruik twee verschillende parameters!) en los op. Een oplossing geeft het snijpunt; geen oplossing betekent kruisend.
\end{werkwijze}

\begin{oefeningen}
\oef{h3-rechten}
\end{oefeningen}

%========================================================================================
\section{Een rechte en een vlak}

Een rechte en een vlak kunnen elkaar in één punt snijden, evenwijdig zijn zonder gemeenschappelijk punt, of de rechte ligt in het vlak. Vul de parametervergelijkingen van de rechte in de vergelijking van het vlak in: dat geeft één vergelijking in de parameter.

\voorbeeld{}{Bereken het snijpunt van de rechte $l\colon (1, 0, 2) + r\vcr{1}{1}{-1}$ met het vlak $\alpha\colon x + 2y + z = 7$.\label{vb:h3-rechtevlak}}{%
\uitleg{Een punt van $l$ is $(1 + r, r, 2 - r)$. Het ligt in $\alpha$ als het de vergelijking klopt.}
\begin{stappenbreed}
(1 + r) + 2r + (2 - r) = 7 \iff 2r = 4 \iff r = 2 & Invullen. \\
\text{snijpunt } (3, 2, 0) & Controle: $3 + 4 + 0 = 7$.
\end{stappenbreed}
Was er een vergelijking zonder oplossing gekomen (zoals $3 = 7$), dan was de rechte evenwijdig met het vlak. Een vergelijking die altijd klopt (zoals $7 = 7$) betekent dat de rechte in het vlak ligt.}

\begin{oefeningen}
\oef{h3-rechtevlak}

\oef{h3-rechtevlak-ligt}
\end{oefeningen}

%========================================================================================
\section{Vlakken}

\begin{theorie}
De vlakken $a_1x + b_1y + c_1z = d_1$ en $a_2x + b_2y + c_2z = d_2$ zijn \textbf{\textcolor{accent}{evenwijdig}} als $(a_1, b_1, c_1)$ en $(a_2, b_2, c_2)$ evenredig zijn; ze \textbf{vallen samen} als ook $d_1$ en $d_2$ in dezelfde verhouding staan. Anders snijden ze elkaar volgens een rechte.
\end{theorie}

\voorbeeld{}{Zijn de vlakken $x + y + z = 3$ en $2x + 2y + 2z = 5$ evenwijdig?\label{vb:h3-vlakken}}{%
\uitleg{Vergelijk de coëfficiënten van $x$, $y$ en $z$, en daarna de rechterleden.}
\begin{stappenbreed}
(2, 2, 2) = 2 \cdot (1, 1, 1) & Evenredig: evenwijdig. \\
5 \neq 2 \cdot 3 & Niet samenvallend.
\end{stappenbreed}}

\begin{oefeningen}
\oef{h3-evenwijdig}
\end{oefeningen}

Drie vlakken snijden elkaar meestal in één punt: de oplossing van een stelsel van drie vergelijkingen met drie onbekenden (zie Matrices). Het kan ook dat ze een rechte gemeen hebben, of geen enkel punt.

\voorbeeld{}{Bereken het gemeenschappelijke punt van de vlakken $x + y + z = 6$, $x - y + 2z = 5$ en $2x + y - z = 1$.\label{vb:h3-drievlakken}}{%
\uitleg{Elimineer telkens een onbekende, zoals bij de methode van Gauss.}
\begin{stappenbreed}
-2y + z = -1 & Tweede min eerste. \\
-y - 3z = -11 & Derde min twee keer de eerste. \\
z = 2y - 1 \Rightarrow -y - 6y + 3 = -11 \iff y = 2 & Invullen. \\
z = 3,\quad x = 6 - 2 - 3 = 1 & Het punt $(1, 2, 3)$.
\end{stappenbreed}}

\begin{oefeningen}
\oef{h3-vlakken}

\oef{h3-drievlakken}
\end{oefeningen}

%========================================================================================
\opdrachtenkop

\begin{oefeningen}
\needspace{16\baselineskip}
\oef{h3-o9}

\needspace{14\baselineskip}
\oef{h3-o12}

\oef{h3-o2}

\oef{h3-o14}

\needspace{14\baselineskip}
\oef{h3-o13}

\oef{h3-o11}

\oef{h3-o1}

\needspace{14\baselineskip}
\oef{h3-o10}

\oef{h3-o3}

\oef{h3-o16}

\oef{h3-o15}

\oef{h3-o19}

\needspace{10\baselineskip}
\oef{h3-o17}

\oef{h3-o6}

\needspace{16\baselineskip}
\oef{h3-o21}

\oef{h3-o18}

\oef{h3-o4}

\oef{h3-o23}

\needspace{14\baselineskip}
\oef{h3-o20}

\oef{h3-o5}

\needspace{16\baselineskip}
\oef{h3-o22}

\needspace{16\baselineskip}
\oef{h3-o24}

\oef{h3-o7}

\oef{h3-o25}

\oef{h3-o8}

\oef{h3-o26}
\end{oefeningen}

%----------------------------------------------------------------------------------------
% KOPPEL: verbanden met andere leerstof
%----------------------------------------------------------------------------------------
\needspace{20\baselineskip}
\opdrachtsectie{Koppel}

\lkr{Matrices, Telproblemen en Logica en bewijzen: kijk na wat al gezien is. Het pakket legt de volgorde niet vast.}

\begin{oefeningen}
\oef{h3-koppel-1}

\oef{h3-koppel-2}

\oef{h3-koppel-3}
\end{oefeningen}

%----------------------------------------------------------------------------------------
\opdrachtsectie{Test jezelf}

\begin{oefeningen}
\oef{h3-t1}

\oef{h3-t5}

\oef{h3-t2}

\oef{h3-t6}

\oef{h3-t3}

\oef{h3-t7}

\oef{h3-t4}
\end{oefeningen}

\herhaalschema{%
  \fv{oef:h3-t1}{vb:h3-snijdend}{oef:h3-snijdend}%
  \fv{oef:h3-t5}{vb:h3-kruisend}{oef:h3-kruisend}%
  \fv{oef:h3-t2}{vb:h3-rechtevlak}{oef:h3-rechtevlak}%
  \fv{oef:h3-t6}{vb:h3-rechtevlak}{oef:h3-rechtevlak-ligt}%
  \fv{oef:h3-t3}{vb:h3-vlakken}{oef:h3-vlakken}%
  \fv{oef:h3-t7}{vb:h3-vlakken}{oef:h3-evenwijdig}%
  \fv{oef:h3-t4}{vb:h3-drievlakken}{oef:h3-drievlakken}}

\ifzelfstudie
\section*{Maak je eigen samenvatting}
Sluit dit boek en schrijf in je schrift: de vier mogelijke liggingen van twee rechten, hoe je ze onderscheidt, de drie mogelijkheden voor een rechte en een vlak, wanneer twee vlakken evenwijdig zijn, en hoe je het snijpunt van drie vlakken berekent. Vergelijk daarna met de kaders.
\fi
